\documentclass[fontsize=11pt]{article}
\usepackage{amsmath}
\usepackage[utf8]{inputenc}
\usepackage[margin=0.75in]{geometry}

\title{CSC111 Project Proposal: TODO FILL IN YOUR PROJECT TITLE HERE}
\author{TODO: FILL IN ALL GROUP MEMBER NAMES HERE}
\date{\today}

\begin{document}
\maketitle

\section*{Problem Description and Research Question}

TODO

\section*{Computational Plan}

\begin{enumerate}
    \item[1.] Parse \verb|courses.html| into \verb|courses.json|, where each course entry stores its course code, raw prerequisite string, and raw exclusion list.

    \item[2.] For each course, parse the raw prerequisite string into a structured, evaluable form:
    \begin{enumerate}
        \item[a)] Extract course codes via regex (\verb|[A-Z]{3}[0-9]{3}[HY][0-9]|).
        \item[b)] Infer logical structure from UofT's formatting conventions: \verb|,| as AND, \verb|/| as OR, parentheses as grouping.
        \item[c)] Classify non-course requirements (credit thresholds, permissions, placement tests) by keyword detection.
    \end{enumerate}

    \item[3.] Implement a custom \verb|CourseTree| class where each node stores a student state (completed courses, total credits). The tree models course selection analogously to \verb|a2_minichess.py|: the current state determines available moves (eligible courses), and selecting a course expands the tree to a new child state.

    \item[4.] Implement core tree methods:
    \begin{enumerate}
        \item[a)] \verb|get_available_courses(state)| --- return all available courses at a current state based on prerequisites/exclusions.
        \item[b)] \verb|append(course)| --- TODO 
        \item[c)] \verb|find_path_to(target)| --- search the tree for a sequence of courses that makes the target course eligible, analogous to searching a game tree for a winning position.
    	\item[d)] \verb|get_credits(optional[Dept])| --- return the number of credits satisfied by the current tree, taking an optional argument to check by department for some requirements.
    \end{enumerate}

    \item[5.] Build an interactive user interface where a student inputs completed courses, views available courses (children of the current node), selects courses to ``take'' (expanding the tree), and queries for paths to a target course.
\end{enumerate}

\section*{References}



% NOTE: LaTeX does have a built-in way of generating references automatically,
% but it's a bit tricky to use so you are allowed to write your references manually
% using a standard academic format like MLA or IEEE.
% See project proposal handout for details.

\end{document}
